\section{The next gate}
\label{sec:tpc46-next-gate}

The preceding theorems split the route beyond TPC--45 into four
logically different tasks.

\subsection{Orbit-frequency regularity of the actual coefficient}

The spectral gap in \cref{sec:tpc46-shift-spectrum} is stable under a
Hilbert vector \(v_{p,r,m}\) that is independent of \(j\), or more
generally under a uniformly band-limited \(j\)-dependence.  The actual
dual coefficient has the form
\begin{equation}
 z_{\mathrm{output\ row},\,j,\,d(m)r}.
 \label{eq:tpc46-next-dual-field}
\end{equation}
No theorem above bounds its orbit bandwidth uniformly over the unit
ball of the dual output space.  A positive continuation would therefore
need either
\begin{enumerate}[label=\textup{(\roman*)},leftmargin=2.4em]
 \item a restriction-stable frequency decomposition of the actual mask
       and output field, with projective mass \(X^{o(1)}\); or
 \item a vector-valued sampling theorem that controls
       \eqref{eq:tpc46-next-dual-field} without first replacing it by a
       scalar smooth profile.
\end{enumerate}
The fixed-column decomposition of TPC--36 supplies the scalar smooth
profile before dualization, not either of these stronger conclusions.

\subsection{Atomic, rather than Cartesian, normalization}

The fixed-shift tensor theorem naturally gives
\begin{equation}
 X^{o(1)}\cD_\otimes
 =X^{o(1)}\sum_h\Delta_h.
 \label{eq:tpc46-next-tensor-scale}
\end{equation}
The physical target instead asks for a bound on the synthesis at one
shift by
\begin{equation}
 X^{1/400+o(1)}\Delta_{h_0}
 \label{eq:tpc46-next-atomic-scale}
\end{equation}
after the actual output identifications.  The examples in
\cref{sec:tpc46-atomic-boundaries} show that no theorem for arbitrary
three coefficient slots can bridge
\eqref{eq:tpc46-next-tensor-scale} and
\eqref{eq:tpc46-next-atomic-scale}.  The next arithmetic input must use
the actual source phases, the projective masks, or cancellation between
residual fibers.  Merely retaining M\"obius symbols is insufficient:
TPC--45 proves that their product is constant on each equal-residual
fiber.

\subsection{Alias unfolding and localization}

The TPC alias family samples the determinant at spacing
\(M_{\rm al}=X^{399/400+o(1)}\).  Its spectrum is the folded sum in
\cref{eq:tpc46-exact-alias-fold}; a gap in the completed consecutive-shift
family does not survive that folding without an additional anti-alias
estimate.  Even an averaged estimate on a consecutive window does not
select \(h_0\) without a maximal or regularity theorem, and the
localization factor in \cref{lem:tpc46-weighted-localization} is sharp.

A useful positive target is therefore a coefficient-specific
inequality of the following schematic form, for \(H_*\ge1\):
\begin{equation}
 \sup_{|h-h_0|\le H_*}
 \frac{\|F_h\|^2}{\Delta_h+\delta_X}
 \ll X^{o(1)}
 \left[
  \frac1{H_*}\sum_{|h-h_0|\le H_*}
  \frac{\|F_h\|^2}{\Delta_h+\delta_X}
 \right],
 \label{eq:tpc46-next-maximal-target}
\end{equation}
where fixed normalization constants in the window average are
inessential.  Here \(\delta_X>0\) is a declared regularizer; one also
needs uniform packet provenance
and a range \(H_*\) that overlaps the analytic averaging scale.  Formula
\eqref{eq:tpc46-next-maximal-target} is a target, not an assumption or a
consequence of the present paper.

\subsection{A reversible continuation}

The next paper may pursue any of the following doors.
\begin{enumerate}[label=\textbf{Door \Alph*.},leftmargin=3.7em]
 \item Prove orbit-frequency decay for one actual projected
       source/mask sector and transfer a nonempty part of the
       shift-spectral gap to that sector.
 \item Establish an atomic-normalized estimate on one additional
       cofactor block, even after averaging a fixed finite set of
       output rows.
 \item Prove an anti-alias or maximal-in-shift theorem connecting the
       completed family to the sparse TPC alias family.
 \item Construct actual-mask coefficients that saturate the
       tensor-to-atomic or fiber frontier.  This would give a sharp
       negative boundary for the present route without asserting that
       the full physical coefficient is large.
\end{enumerate}
Success at any door is a genuine advance.  Failure at one does not
invalidate the spectral gap, tensor closure, or sharp boundary results
proved here.  None of these doors may be relabelled as a fixed-shift
prime-pair asymptotic or a breach of the parity barrier.
